\documentclass{article}
\usepackage[utf8]{inputenc}
\usepackage{amsmath}
\usepackage{amssymb}

\title{核心数学模型}
\date{}

\usepackage[UTF8]{ctex}
\begin{document}

\maketitle

\section*{在需求不确定的情况下，云服务器部署策略}

来源：第3.2节，PDF第10-14页，(1)-(18节)；第5.2节，PDF第23页。
这是一个有限情景下的CVaR版本，属于完整的双阶段云服务器部署模型。 它将这两个阶段视为一个整体问题来处理。

\paragraph{集合与数据。}
$\mathcal L$表示数据中心，$\mathcal S$表示供应商，
$\mathcal D^1$表示已知的集群请求，而$\mathcal D^2_\omega$则代表在场景$\omega\in\Omega$中新出现的请求。这些场景具有相同的权重，与源代码的抽样公式和保留的输入数据一致。
$\mathcal T=\{0,\ldots,T\}$、$\mathcal T_1=\{0,\ldots,T_1\}$和$\mathcal T_2=\{T_1+1,\ldots,T\}$。现在，只有$\mathcal T_2$中的行是可用。$\mathcal L(d),\mathcal S(d)$及其以场景索引表示的版本是兼容的数据中心和供应商。$\mathcal H$包含对$(p,t)$的配对，其中$p\subseteq\mathcal L$是一个可能重叠的层次化吞吐量组。$c_{d,t},u_d,h_{d,s}$表示非负的早期/晚期对接、未满足的需求和运输成本；上标$\omega$表示场景数据。$b_\ell$是行构建成本，$\zeta_\ell$是总行容量，而$\rho_{\ell,1},\sigma_{s,1},
\delta_{p,t,1}$则是初始空闲行、供应商库存和吞吐量。容量数据是非负整数。不兼容的变量在求和时被视为零。

\paragraph{变量。}
$x_\ell$表示新增的行数。二进制表示法。  
$z_{d,\ell,t}$用于记录集群请求，而$w_{d,s}$则用于选择其供应商。  
第二阶段的相关数值以$\omega$表示。剩余资源$\rho_2,\sigma_2,\delta_2$由所有场景共享，无法预测其具体数值。$q_\omega$表示完整的非负场景成本。$\eta\in\mathbb R$、$e_\omega\ge0$和$\theta\in\mathbb R$则分别定义了在置信水平$\alpha\in[0,1)$下的CVaR值。论文中的尾部质量表示方式为$1-\alpha$。

\paragraph{单一确定性等价。}
\begin{equation}
\min
\sum_{d\in \mathcal{D}^1}
\left[
\sum_{\ell\in \mathcal{L}(d)} \sum_{t\in \mathcal{T}} c_{d,t} z_{d,\ell,t}
+ u_d \left(1 - \sum_{\ell\in \mathcal{L}(d)} \sum_{t\in \mathcal{T}} z_{d,\ell,t}\right)
+ \sum_{s\in \mathcal{S}(d)} h_{d,s} w_{d,s}
\right]
+ \sum_{\ell\in \mathcal{L}} b_\ell x_\ell
+ \theta
\tag{1}
\end{equation}

\begin{equation}
\sum_{\ell\in \mathcal{L}(d)} \sum_{t\in \mathcal{T}} z_{d,\ell,t} \le 1,
\qquad \forall d \in \mathcal{D}^1
\tag{2}
\end{equation}

\begin{equation}
\sum_{\ell\in \mathcal{L}(d)} \sum_{t\in \mathcal{T}} z_{d,\ell,t}
=
\sum_{s\in \mathcal{S}(d)} w_{d,s},
\qquad \forall d \in \mathcal{D}^1
\tag{3}
\end{equation}

\begin{equation}
\delta_{p,t,1}
\ge
\sum_{d\in \mathcal{D}^1} \sum_{\ell\in p \cap \mathcal{L}(d)} z_{d,\ell,t},
\qquad \forall (p,t) \in \mathcal{H}:\ t \in \mathcal{T}_1
\tag{4}
\end{equation}

\begin{equation}
\delta_{p,t,2}
=
\delta_{p,t,1}
- \sum_{d\in \mathcal{D}^1} \sum_{\ell\in p \cap \mathcal{L}(d)} z_{d,\ell,t},
\qquad \forall (p,t) \in \mathcal{H}:\ t \in \mathcal{T}_2
\tag{5}
\end{equation}

\begin{equation}
\rho_{\ell,1}
\ge
\sum_{d\in \mathcal{D}^1:\ell\in \mathcal{L}(d)} \sum_{t\in \mathcal{T}_1} z_{d,\ell,t},
\qquad \forall \ell \in \mathcal{L}
\tag{6}
\end{equation}

\begin{equation}
\rho_{\ell,2}
=
\rho_{\ell,1} + x_\ell
- \sum_{d\in \mathcal{D}^1:\ell\in \mathcal{L}(d)} \sum_{t\in \mathcal{T}} z_{d,\ell,t},
\qquad \forall \ell \in \mathcal{L}
\tag{7}
\end{equation}

\begin{equation}
\zeta_\ell \ge \rho_{\ell,1} + x_\ell,
\qquad \forall \ell \in \mathcal{L}
\tag{8}
\end{equation}

\begin{equation}
\sigma_{s,2}
=
\sigma_{s,1} - \sum_{d\in \mathcal{D}^1:s\in\mathcal S(d)} w_{d,s},
\qquad \forall s \in \mathcal{S}
\tag{9}
\end{equation}

\begin{equation}
\rho_{\ell,2},\ \delta_{p,t,2},\ \sigma_{s,2} \ge 0
\tag{10}
\end{equation}

\begin{equation}
z_{d,\ell,t},\ w_{d,s} \in \{0,1\},
\qquad
x_\ell \in \mathbb{Z}_+
\tag{11}
\end{equation}

\begin{equation}
q_\omega
=
\sum_{d\in \mathcal{D}^2_\omega}
\left[
\sum_{\ell\in \mathcal{L}_\omega(d)} \sum_{t\in \mathcal{T}_2} c^\omega_{d,t} z^\omega_{d,\ell,t}
+ u^\omega_d \left(1 - \sum_{\ell\in \mathcal{L}_\omega(d)} \sum_{t\in \mathcal{T}_2} z^\omega_{d,\ell,t}\right)
+ \sum_{s\in \mathcal{S}_\omega(d)} h^\omega_{d,s} w^\omega_{d,s}
\right],
\qquad \forall \omega \in \Omega
\tag{12}
\end{equation}

\begin{equation}
q_\omega \ge 0,
\qquad \forall \omega \in \Omega
\tag{12a}
\end{equation}

\begin{equation}
\sum_{\ell\in \mathcal{L}_\omega(d)} \sum_{t\in \mathcal{T}_2} z^\omega_{d,\ell,t} \le 1,
\qquad \forall d \in \mathcal{D}^2_\omega,\ \forall \omega \in \Omega
\tag{13}
\end{equation}

\begin{equation}
\sum_{\ell\in \mathcal{L}_\omega(d)} \sum_{t\in \mathcal{T}_2} z^\omega_{d,\ell,t}
=
\sum_{s\in \mathcal{S}_\omega(d)} w^\omega_{d,s},
\qquad \forall d \in \mathcal{D}^2_\omega,\ \forall \omega \in \Omega
\tag{14}
\end{equation}

\begin{equation}
\delta_{p,t,2}
\ge
\sum_{d\in \mathcal{D}^2_\omega} \sum_{\ell\in p \cap \mathcal{L}_\omega(d)} z^\omega_{d,\ell,t},
\qquad \forall (p,t) \in \mathcal{H}:\ t \in \mathcal{T}_2,\ \forall \omega \in \Omega
\tag{15}
\end{equation}

\begin{equation}
\rho_{\ell,2}
\ge
\sum_{d\in \mathcal{D}^2_\omega:\ell\in \mathcal{L}_\omega(d)} \sum_{t\in \mathcal{T}_2} z^\omega_{d,\ell,t},
\qquad \forall \ell \in \mathcal{L},\ \forall \omega \in \Omega
\tag{16}
\end{equation}

\begin{equation}
\sigma_{s,2}
\ge
\sum_{d\in \mathcal{D}^2_\omega:\ s\in \mathcal{S}_\omega(d)} w^\omega_{d,s},
\qquad \forall s \in \mathcal{S},\ \forall \omega \in \Omega
\tag{17}
\end{equation}

\begin{equation}
z^\omega_{d,\ell,t},\ w^\omega_{d,s} \in \{0,1\},
\qquad \forall d,\ell,t,s,\omega
\tag{18}
\end{equation}


\begin{align}
\theta&=\eta+\frac{1}{(1-\alpha)|\Omega|}\sum_{\omega\in\Omega}e_\omega,\tag{19}\\
e_\omega&\ge q_\omega-\eta,\quad e_\omega\ge0 &&\omega\in\Omega,\tag{20}\\
\eta,\theta&\in\mathbb R.\tag{21}
\end{align}
第一阶段的目标（1）需要被最小化一次，这一过程需要同时考虑所有变量和约束条件。在每个情景中，所有缺失的常数都会被保留下来。
\end{document}
